\section{Overview}%
\label{sec:chap5Overview}
\autoBookmark{sec}{Overview}
This chapter closes the DSRM cycle. It revisits each research objective against the mapping of Table~\ref{tab:dsrm_mapping}, states what is new in the work and what it contributes to knowledge, sets out the limitations, and identifies where the work should go next. Each conclusion is stated together with the scope within which it holds, because several findings of Chapter~\ref{chap:ResultsAndDiscussion} apply to recordings represented in training but not to new acquisitions or to a new sensing system.

\section{Revisiting the Research Objectives}
\label{sec:chap5Objectives}
\autoBookmark{sec}{Revisiting the Research Objectives}
An objective is regarded as met when the design, implementation, comparison, or measurement it specifies was carried out as planned and the corresponding research question can be answered from the results, whether or not the answer is the one anticipated. Where an answer differs from expectation, it is reported as measured.

\subsection{RO1: An Interpretable Pipeline Under Leakage-Controlled Evaluation}
\label{sec:chap5RO1}
\autoBookmark{subsec}{RO1: An Interpretable Pipeline Under Leakage-Controlled Evaluation}
An interpretable, feature-driven pipeline was designed on the Tomasov dataset and applied unchanged to the Cao dataset. It combines a stateless three-stage preprocessing chain, a single-pass extractor of $1{,}002$ named descriptors across the time, frequency, time--frequency, and spatial domains, and seven lightweight classifiers, and it needs no gradient-based training on the edge device. Its evaluation is part of the design. Auditing the released protocols showed that every test window of a random Tomasov split shares raw samples with a training window and that 423 of the 441 Cao sessions span both published folders. Blocked, guard-banded partitions for Tomasov and session-grouped partitions for Cao reduce both measures to zero, and their disjointness is asserted in code before any model is trained. The protocol-attribution experiments show why this matters: changing only the partition moved an otherwise identical CNN by $26.27$ accuracy points and five recent architectures by $6.5$ to $18.4$. The selected LightGBM classifies a window in under $0.2$\,ms on embedded hardware, but the lightweight claim applies to the classifier, not to the whole pipeline: extraction takes $144.59$\,ms per window on Tomasov and $122.54$\,ms on Cao, so Tomasov needs two cores to keep up with its arrival rate. The verification covers the partitions rather than the feature code, and a later audit of that code found the numerical residue reported in \refss{sec:chap5LimResidue}. RO1 is met, with that qualification.

\subsection{RO2: Multi-Domain Fusion, Feature Selection, and Class Balancing}
\label{sec:chap5RO2}
\autoBookmark{subsec}{RO2: Multi-Domain Fusion, Feature Selection, and Class Balancing}
Fusing the four domains outperformed every individual domain on both datasets, by $6.26$ macro-$F_1$ points on Tomasov and $3.46$ on Cao, so the domains carry complementary rather than redundant information. The domain ordering agreed across the two systems at $\rho = 0.90$, which points to a property of the representation rather than of one campus fibre, although the agreement weakens once the residue of \refss{sec:chap5LimResidue} is removed.

The eight-method consensus, with every selector confined to training data and refitted inside each cross-validation fold, reduced the feature space by about $95\%$, to 52 features on Tomasov and 53 on Cao. Two findings qualify this result. The reduction is not free: it costs $8.39$ and $5.81$ macro-$F_1$ points, far more than the sub-point penalty reported in earlier single-dataset work, so an intermediate operating point on the reduction curve is preferable unless inference cost or interpretability dominates. And the six balancing strategies separated only modestly on Tomasov and not at all on Cao when compared over independent held-out acquisitions, so the expectation that cost-sensitive weighting would be consistently superior was not supported. RO2 is met, with both qualifications reported as measured.

\subsection{RO3: Comparative Performance, Cost, and Generalisation}
\label{sec:chap5RO3}
\autoBookmark{subsec}{RO3: Comparative Performance, Cost, and Generalisation}
Seven classifiers were benchmarked across six balancing strategies on both datasets. The three gradient-boosting learners led on both and lay within $0.62$ and $0.84$ macro-$F_1$ points of one another, so LightGBM, which validation selected in every configuration, is best regarded as the representative of that tier rather than as a learner shown to beat the other two. It reached $90.71\%$ accuracy and $81.98$ macro-$F_1$ on Tomasov and $99.06\%$ and $98.90$ on Cao, and its lead over random forest and logistic regression held on all five independently generated partitions.

Against the CNN published with each dataset, reimplemented from the released code or the paper's specification and trained on identical partitions, the pipeline led by $16.43$ and $8.01$ macro-$F_1$ points; the Tomasov reconstruction exceeds its own published figure under the original protocol, so it is not a weakened comparator. Against five architectures published from 2024 onwards, trained under the same partitions, the pipeline led the strongest neural baseline by $11.88$ accuracy points on Tomasov and $1.11$ on Cao. With the residue of \refss{sec:chap5LimResidue} removed, the Tomasov lead narrows to $1.69$ points and S2-DyGNN moves ahead on macro-$F_1$. The neural baselines, which consume raw windows, are nevertheless between twenty and nine hundred times cheaper online, so the pipeline's advantages are interpretability and training cost rather than latency.

Generalisation was tested under progressively stricter protocols. Accuracy averaged $63.85\%$ (SD $29.31$) on unseen recordings of the four Tomasov classes recorded more than once, and fell from $99.06\%$ to $86.81\%$ on unseen Cao acquisition dates. Mapping the two taxonomies onto four shared event types showed that these events are recognised well within each system ($73.02\%$ and $95.93\%$ accuracy from non-spatial features) but that a model trained on one system scores below chance on the other ($14.41\%$ and $19.12\%$). Repeating the leakage-controlled partition with five seeds gave a Tomasov standard deviation of $4.42$ macro-$F_1$ points, so every headline figure on that dataset carries a substantial partition uncertainty. RO3 is met: the comparisons and cost measurements were made as planned, and the generalisation limits are reported as found.

\section{Contributions to Knowledge}
\label{sec:chap5Contributions}
\autoBookmark{sec}{Contributions to Knowledge}
The novelty of this thesis does not lie in a new learning algorithm. Gradient boosting, wavelet descriptors, SMOTE, and blocked cross-validation are all established. What is new, seen from information technology and computer science, is a way of building and measuring a DAS classifier in which the validity of the evaluation is a tested property of the software. Provenance is recorded before any signal is read, a partition that fails its checks stops the run, each claim about a model is tied to a controlled experiment that changes one factor, and cost is measured for the whole pipeline rather than for the classifier alone.

In the knowledge-contribution framework of Gregor and Hevner~\cite{gregor2013positioning}, the evaluation protocol is an \emph{exaptation}: blocked and group-aware evaluation are known solutions from other fields, carried here into a field that had not applied them. The attribution design, the recording-identity probe, the invariance audit, and the feature pipeline are \emph{improvements}, new solutions to a known problem. No claim of invention is made. Table~\ref{tab:novelty_evidence} sets each element introduced in \refss{sec:ThesisNovelty} beside the evidence for it.

\begin{table}[!htbp]
	\centering
	\caption[Novel elements of the thesis and the evidence for each]{Novel elements of the thesis, the kind of knowledge contribution each makes in the framework of Gregor and Hevner~\cite{gregor2013positioning}, and the measured evidence for each.}
	\label{tab:novelty_evidence}
	\footnotesize
	\setlength{\tabcolsep}{4pt}
	\renewcommand{\arraystretch}{1.2}
	\begin{tabularx}{\textwidth}{@{}>{\raggedright\arraybackslash}p{4.1cm} >{\raggedright\arraybackslash}p{1.9cm} >{\raggedright\arraybackslash}X@{}}
		\toprule
		\textbf{Novel element} & \textbf{Type} & \textbf{Evidence in this thesis}\\
		\midrule
		Evaluation validity asserted in code, with provenance manifests & Exaptation & Test windows sharing raw samples with training fall from $100\%$ to $0\%$ on Tomasov; sessions spanning the split fall from 423 of 441 to none on Cao.\\
		\addlinespace[3pt]
		Blocked partitioning with guard bands derived from read footprints & Exaptation & All nine classes stay in every split, including the five single-recording classes, at a cost of $9.00\%$ of the annotated windows.\\
		\addlinespace[3pt]
		Protocol-attribution experiments & Improvement & Changing only the partition moves an identical CNN by $26.27$ accuracy points and five recent architectures by $6.48$ to $18.38$.\\
		\addlinespace[3pt]
		Recording-identity probe & Improvement & Spatial descriptors identify the source recording $11.1$ points beyond the next domain, yet removing them costs $19.08$ accuracy points on unseen recordings.\\
		\addlinespace[3pt]
		Four-domain representation with unanimous consensus selection & Improvement & Fusion leads the best single domain by $6.26$ and $3.46$ macro-$F_1$ points; consensus removes about $95\%$ of the features at a cost of $8.39$ and $5.81$.\\
		\addlinespace[3pt]
		Whole-pipeline cost accounting on embedded hardware & Improvement & On Tomasov, extraction takes $144.59$\,ms and inference $0.17$\,ms against a $102.4$\,ms window interval, so extraction sets the real-time budget.\\
		\addlinespace[3pt]
		Mechanism-based taxonomy mapping and cross-system transfer test & Exaptation & Four shared events are recognised at $73.02\%$ and $95.93\%$ accuracy within each system but at $14.41\%$ and $19.12\%$ across systems, below the $25\%$ chance level.\\
		\addlinespace[3pt]
		Invariance audit of the feature code & Improvement & Thirty-three of the $1{,}002$ descriptors respond to changes that should not reach them, and the residue accounts for about ten Tomasov accuracy points (\refss{sec:chap5LimResidue}).\\
		\bottomrule
	\end{tabularx}
\end{table}

These elements yield six findings that matter for DAS research more broadly.

The first is that the evaluation protocol, not the architecture alone, governs the headline results on this benchmark. Holding a model fixed and changing only the partition altered its accuracy by $26.27$ points, and the same intervention changed five recent architectures by $6.5$ to $18.4$ points. Published accuracies obtained under random or chronological partitions of overlapping windows are therefore not directly comparable and should not be read as model rankings.

The second is that leakage substantially inflated minority-class performance, while the feature representation remained the limiting factor. Once contamination was removed, the \emph{openclose} class fell from a published $F_1$ of $0.72$ to $0.044$ for the selected model, and none of the seven classifiers exceeded $0.19$. The class is not unlearnable from the data: on the same partition and the same 198 test windows, S2-DyGNN reaches $0.625$, GCN $0.433$, and GAT $0.278$, each above every feature-based classifier. It is the pooled descriptors used here, whose 61 spatial features aggregate across the channel axis, that cannot recover it. Because \emph{openclose} corresponds to unauthorised manhole access, the correction matters operationally, and it identifies a case in which pooled descriptors should be replaced or supplemented by a model that keeps the channel axis.

The third is that a feature set can be informative and confounded at the same time. In the recording-identity probe, spatial descriptors identified the source recording with a margin $11.1$ points larger than that of the next domain, yet leave-one-recording-out evaluation showed that removing them reduced cross-recording accuracy by $19.08$ points. Both findings hold, and they cannot be separated in a dataset in which five of the nine classes come from a single recording. The methodological contribution is the probe itself, a diagnostic that quantifies the confound rather than presuming it absent.

The fourth is that statistics must be computed over independent acquisitions. The Friedman test, applied with balancing strategies as blocks, returned a near-maximal statistic that reflected the dependence between blocks rather than an effect, and it was withdrawn. Repeating the partition with five seeds gave a Tomasov standard deviation of $4.42$ macro-$F_1$ points, and bootstrapping over recordings rather than windows widened the confidence interval by an order of magnitude. Any headline figure on a dataset with few independent recordings must carry that uncertainty.

The fifth is that a method can generalise across sensing systems while a fitted model does not. The same pipeline configuration, untuned, recognised the shared event types within each of two unrelated systems, yet a model fitted on one system and applied to the other, on four carefully mapped event types, scored below chance in both directions. Cross-system claims in DAS should therefore state which of the two is meant, and the mapping between the Tomasov and Cao taxonomies gives future adaptation methods a concrete benchmark on which to show progress.

The sixth concerns leakage below the level of the partition. A constant of $10^{-12}$ added to a denominator for numerical safety left about $10^{-14}$ of each channel's raw amplitude in the features, and tree ensembles, which compare values rather than weigh them, could split on it (\refss{sec:chap5LimResidue}). No partition check can see such a path, because it runs through the feature code rather than through the split. A simple invariance audit can: multiply each raw channel by its own gain and test which features move. Any pipeline that claims amplitude normalisation, in DAS or elsewhere, can run this check on a few windows per recording, and the audit is released with this thesis.

\section{Limitations}
\label{sec:chap5Limitations}
\autoBookmark{sec}{Limitations}
The conclusions above hold within limits, set out here as threats to validity. \refss{sec:chap5LimConfound} and \refss{sec:chap5LimResidue} concern internal validity, whether measured performance reflects the event rather than the recording, and they bear most directly on the headline figures. \refss{sec:chap5LimMethod} concerns the strength of the comparisons, and \refss{sec:chap5LimScope} concerns external validity, how far the findings extend beyond these datasets and this offline setting.

\subsection{The Recording Confound Cannot Be Resolved with These Datasets}
\label{sec:chap5LimConfound}
\autoBookmark{subsec}{The Recording Confound Cannot Be Resolved with These Datasets}
The central structural limitation lies in the data. Because five Tomasov classes each come from a single recording, recording identity and class identity are indistinguishable for them, and no partition can separate their effects. Two operating points are therefore reported, but definitive attribution needs a dataset with several independent recordings per class. The Tomasov probe also rests on limited evidence: after degenerate cases are excluded, only one class remains suitable for analysis, so that measurement should be regarded as indicative. Five probes remain valid on Cao, which gives stronger evidence of a dataset-wide campaign confound. The Tomasov labels themselves were drawn by hand on band-limited power maps, and their noise was not estimated.

\subsection{A Numerical Residue in Amplitude Normalisation}
\label{sec:chap5LimResidue}
\autoBookmark{subsec}{A Numerical Residue in Amplitude Normalisation}
The z-score of \refss{sec:chap3Normalisation} was implemented with the customary stabilising constant, dividing by $\sigma_{\mathcal{H}} + \epsilon$ with $\epsilon = 10^{-12}$. After the experiments of Chapter~\ref{chap:ResultsAndDiscussion} had been completed, an invariance audit tested the feature code directly: each raw channel was multiplied by its own gain between $0.25$ and $4$, each raw sample was perturbed by about one part in $10^{12}$, and every descriptor was compared before and after. Per-channel z-scoring should make all $1{,}002$ descriptors indifferent to such changes. Thirty-three were not: thirteen carried each channel's raw amplitude through the constant, and the rest followed floating-point rounding. A channel with standard deviation $\sigma_c$ ends up with energy $T(1 - 2\epsilon/\sigma_c)$ rather than exactly $T$, an imprint of about $10^{-14}$. That is far below any physical scale and close to the limit of double-precision arithmetic~\cite{goldberg1991floating}, but a decision tree compares values rather than weighing them and can place a split between any two distinct numbers. The residue therefore reopens, in numerical form, the fingerprint of gain and coupling that normalisation is meant to close. The partition checks cannot detect it, since it involves no shared samples.

Chapter~\ref{chap:ResultsAndDiscussion} reports the experiments as executed, with the constant in place. Table~\ref{tab:residue_effect} sets those figures beside the same measurements with the chain made exact: dividing by $\sigma_{\mathcal{H}}$ alone, writing the descriptors that are constant by construction as exact constants, and flooring the spectrum at $-120$\,dB.

\begin{table}[!htbp]
	\centering
	\caption{Effect of the numerical residue: the figures reported in Chapter~\ref{chap:ResultsAndDiscussion}, with the constant in place, beside the same measurements with the chain made exact.}
	\label{tab:residue_effect}
	\small
	\setlength{\tabcolsep}{4pt}
	\begin{tabular}{@{}lrr@{}}
		\toprule
		\textbf{Measure} & \textbf{As reported} & \textbf{Exact chain}\\
		\midrule
		\multicolumn{3}{@{}l}{\textit{Tomasov}}\\[2pt]
		Accuracy of the selected model (\%)                    & 90.71 & 80.52\\
		Macro-$F_1$ of the selected model                      & 81.98 & 76.33\\
		\emph{Openclose} $F_1$ of the selected model           & 0.044 & 0.263\\
		Macro-$F_1$ without the spatial block                  & 62.13 & 62.07\\
		Macro-$F_1$ lost on removing the spatial block         & 19.67 & 14.84\\
		Spatial recording identification above chance (points) & 12.9  & 8.8\\
		Leave-one-recording-out accuracy (\%)                  & 63.85 & 53.81\\
		Macro-$F_1$ over five partitions (mean $\pm$ SD)       & $80.47 \pm 4.42$ & $73.04 \pm 3.22$\\
		\midrule
		\multicolumn{3}{@{}l}{\textit{Cao}}\\[2pt]
		Accuracy of the selected model (\%)                    & 99.06 & 98.50\\
		Macro-$F_1$ of the selected model                      & 98.90 & 98.18\\
		Leave-one-date-out accuracy (\%)                       & 86.81 & 85.98\\
		\midrule
		\multicolumn{3}{@{}l}{\textit{Across the two systems}}\\[2pt]
		Transfer accuracy, Cao to Tomasov (\%)                 & 14.41 & 15.76\\
		Transfer accuracy, Tomasov to Cao (\%)                 & 19.12 & 20.25\\
		Rank correlation of the feature domains ($\rho$)       & 0.90  & 0.50\\
		\bottomrule
	\end{tabular}
\end{table}

The effect is large on Tomasov and small on Cao. Tomasov accuracy falls by $10.19$ points to $80.52\%$ and macro-$F_1$ by $5.65$ to $76.33$, with XGBoost rather than LightGBM as the validated choice, which is no surprise given that the boosted learners cannot be separated (\refss{sec:resAllModels}). Cao moves by $0.56$ and $0.72$. The Tomasov headline of Chapter~\ref{chap:ResultsAndDiscussion} should therefore be read as an upper bound, with about ten of its accuracy points owed to the residue rather than to the events.

Several conclusions change in degree. The lead over the published Tomasov CNN falls from $16.43$ to $10.78$ macro-$F_1$ points. The lead over the strongest neural baseline falls from $11.88$ to $1.69$ accuracy points, and on macro-$F_1$ the $78.54$ of S2-DyGNN moves ahead of $76.33$. The residue sits in the channel-energy descriptors, so the contribution of the spatial domain and its recording-identity margin both shrink, and the fall of that domain on Cao lowers the agreement between the two systems from $\rho = 0.90$ to $0.50$. Cross-recording accuracy drops by a further ten points.

Others hold. Neither the leakage audit, nor the partitions, nor any neural baseline reads these features, so the protocol-attribution results of \refs{sec:resDiscussion} are untouched. Fusion still leads the best single domain on both datasets, by $6.22$ and $5.84$ macro-$F_1$ points. The spatial domain remains the strongest single domain on Tomasov, and the conservative no-spatial figure of \refss{sec:resConfoundVerdict} barely moves. \emph{Openclose} stays below all three graph models, and cross-system transfer stays below chance.

\subsection{Methodological Scope}
\label{sec:chap5LimMethod}
\autoBookmark{subsec}{Methodological Scope}
Several comparisons rest on single runs or on choices that were fixed rather than optimised. Each CNN and neural baseline was trained once, at one seed. On a class of 198 test windows a single run is a weak estimate, so the \emph{openclose} values of Table~\ref{tab:sota} show that the class carries signal a channel-level model can use, not what these models reach reliably, and the conclusion rests on the pattern across the three graph models rather than on the $0.625$ of S2-DyGNN alone. The classical learners share one hyperparameter configuration, which keeps the comparison fair between them but does not establish the best each could reach. In the consensus ensemble, mRMR scores 400 features while the retention rule keeps 501, so 101 of its votes fall to tie-breaking. Some features therefore leave the unanimous subset for reasons unconnected with their merit, and the reported subset is probably smaller, and its cost larger, than a full ranking would give.

The baselines are reconstructions. Where a recent architecture's publication left a detail unstated, the choice made is recorded rather than tuned, so those baselines are faithful reconstructions rather than the authors' own code. Because librosa was not usable on the embedded platform, the MFCC branch of DASViT1D uses a SciPy implementation that yields 980 values rather than the $1{,}300$ of the librosa configuration, so its results are specific to that front end. The RDFT redundancy factor for the Tomasov CNN was inferred from two stated quantities rather than read directly, and only the RDFT front end was evaluated, not the MFCC and DWT variants the same publication reports. The re-cut partition used for protocol attribution reconstructs the published protocol rather than reproducing it, and single-run CNN accuracies on Tomasov vary by a few tenths with row ordering alone. The Cao sampling rate is a nominal value because the release does not identify which fibre span produced each record, and the Cao CNN was not evaluated under leave-one-date-out, so the capacity-versus-leakage observation is established for within-campaign evaluation only. Finally, the verification code is written for the metadata of these two datasets, and applying it elsewhere needs engineering effort.

\subsection{Deployment Scope}
\label{sec:chap5LimScope}
\autoBookmark{subsec}{Deployment Scope}
Within the scope of these datasets, a fitted model should be regarded as valid only for the acquisition campaign and the interrogator that produced its training data. The two datasets differ in almost every respect: Tomasov has nine classes from 16 phase-sensitive recordings over $1{,}700$ loci, and Cao six classes from 441 intensity-based sessions over 12 channels, with different sampling, geometry, labels, and sample units. Applying one pipeline to both shows that the method carries over, but the cross-system transfer test shows that a fitted model does not. That test covers only the four event types with a one-to-one counterpart, and those correspondences are approximate: construction machinery and manual digging, or a fence strike and a knock on the ground, excite the fibre differently. Robustness to weather, seasonal variation, soil properties, burial depth, source direction relative to the fibre, and other interrogator hardware was not evaluated.

The evaluation was also offline. Every result comes from stored recordings; the pipeline has not been connected to a live interrogator, and its energy draw was not measured. The real-time budget of Table~\ref{tab:efficiency} concerns one stream of windows at one locus, whereas an interrogator reports $1{,}700$ loci. The Tomasov test set is decimated in the same proportion as the training set, so its accuracy and error rates do not reflect the prevalence of routine activity on a live fibre, where \emph{regular} would dominate and false alarms would weigh more heavily. Interpretability, finally, is argued from named descriptors and attribution scores; it has not been tested with the operators who would act on the alerts.

\section{Future Directions}
\label{sec:chap5Future}
\autoBookmark{sec}{Future Directions}
The limitations point to seven directions, presented in order of expected value (Table~\ref{tab:limits_future}). The first three strengthen the evidence and make it reusable, the next two improve the model, and the last two move the artefact towards use. Each builds directly on an artefact or finding of this thesis rather than requiring it to be rebuilt.

\begin{table}[!htbp]
	\centering
	\caption{From limitations to future work.}
	\label{tab:limits_future}
	\footnotesize
	\setlength{\tabcolsep}{4pt}
	\renewcommand{\arraystretch}{1.2}
	\begin{tabularx}{\textwidth}{@{}>{\raggedright\arraybackslash}X >{\raggedright\arraybackslash}X >{\raggedright\arraybackslash}p{1.6cm}@{}}
		\toprule
		\textbf{Limitation} & \textbf{Direction} & \textbf{Subsection}\\
		\midrule
		Numerical residue in the reported figures; single training runs for the baselines & Re-baseline with exact normalisation and repeated training & \ref{sec:chap5FutureBaseline}\\
		\addlinespace[3pt]
		Recording and class cannot be separated for five Tomasov classes & Collect a dataset with several recordings per class & \ref{sec:chap5FutureDataset}\\
		\addlinespace[3pt]
		Verification code written for two datasets & Package the controls as reusable software & \ref{sec:chap5FutureTooling}\\
		\addlinespace[3pt]
		Pooled descriptors cannot recover \emph{openclose} & Combine them with a channel-level model & \ref{sec:chap5FutureChannelModel}\\
		\addlinespace[3pt]
		Fitted models fail on unseen recordings, campaigns, and systems & Pursue domain adaptation & \ref{sec:chap5FutureAdaptation}\\
		\addlinespace[3pt]
		Extraction exceeds the Tomasov window interval; evaluation is offline and at one locus & Reduce extraction cost and move to streaming deployment & \ref{sec:chap5FutureCost}\\
		\addlinespace[3pt]
		Interpretability argued but not tested with users & Evaluate interpretability with operators & \ref{sec:chap5FutureOperators}\\
		\bottomrule
	\end{tabularx}
\end{table}

\subsection{Re-Baseline the Benchmark}
\label{sec:chap5FutureBaseline}
\autoBookmark{subsec}{Re-Baseline the Benchmark}
The first step is inexpensive. Every table of Chapter~\ref{chap:ResultsAndDiscussion} should be regenerated with the exact chain of Table~\ref{tab:residue_effect}, and the CNNs and neural baselines should be trained over several seeds, so that per-class values, \emph{openclose} above all, carry an interval rather than a point. The invariance audit should then run as a standing gate beside the partition checks, so that a later change to the feature code cannot reopen the same path unnoticed.

\subsection{Collect a Multi-Recording Dataset}
\label{sec:chap5FutureDataset}
\autoBookmark{subsec}{Collect a Multi-Recording Dataset}
Repeating each event class across separate fibre sections, days, and coupling conditions would allow the recording confound to be resolved rather than merely measured, and would permit leave-one-recording-out evaluation for the full taxonomy. This is the most valuable new evidence for this benchmark. Such a dataset would also let the recording-identity probe be applied to every class rather than to the single Tomasov class that survives its exclusion criteria, and the partition and verification procedures developed here could be applied to it unchanged.

\subsection{Package the Leakage Controls as Reusable Software}
\label{sec:chap5FutureTooling}
\autoBookmark{subsec}{Package the Leakage Controls as Reusable Software}
The manifests, the partition assertions, the attribution harness, and the invariance audit are written for two datasets. Generalised into an installable library, driven by a declarative description of a dataset's dependence groups, they could run as an automated test in any DAS project, much as continuous integration runs unit tests, and fail a build when a split shares samples. Frozen public manifests could then support a shared leaderboard on which results are comparable because their protocol is the same. Software-engineering work on machine-learning systems has argued for this kind of automated readiness testing~\cite{sculley2015debt,breck2017mltest}, and DAS benchmarks are an open place to apply it.

\subsection{Combine Pooled Descriptors with a Channel-Level Model}
\label{sec:chap5FutureChannelModel}
\autoBookmark{subsec}{Combine Pooled Descriptors with a Channel-Level Model}
The \emph{openclose} result shows that graph models operating over the channel axis recover information that aggregate spatial descriptors lose, while the feature pipeline is stronger on the other classes. A natural extension would add a graph model's channel-level embedding to the engineered feature space, or route low-confidence windows to the graph model. Because the two approaches fail in different places, such a hybrid could keep the interpretability of the descriptors for the classes they handle well while delegating the channel-dependent class to a model that preserves the channel axis.

\subsection{Pursue Domain Adaptation Rather than Further Feature Engineering}
\label{sec:chap5FutureAdaptation}
\autoBookmark{subsec}{Pursue Domain Adaptation Rather than Further Feature Engineering}
A twelve-point fall from within-campaign to cross-campaign accuracy on Cao, a twenty-seven-point fall to unseen recordings on Tomasov, and below-chance transfer between the two systems all point to distribution shift as the binding constraint on deployment. Transfer learning, domain-adversarial training, explicit calibration, and open-set detection of unseen events are the appropriate responses. The leave-one-date-out and leave-one-recording-out protocols of this thesis, together with the four-class mapping between the Tomasov and Cao taxonomies, provide an evaluation framework against which such methods can be measured without reintroducing leakage.

\subsection{Reduce Extraction Cost and Move to Streaming Deployment}
\label{sec:chap5FutureCost}
\autoBookmark{subsec}{Reduce Extraction Cost and Move to Streaming Deployment}
Feature extraction, not inference, sets the real-time budget, and the time--frequency domain accounts for about $60$\,ms per window. GPU-accelerated wavelet-packet decomposition, or removing texture descriptors that the consensus subsets rarely retain, would improve deployability more than further classifier tuning. Because the consensus subsets keep only 52 to 53 of the $1{,}002$ features, extracting only those descriptors would cut the per-window cost further, at the accuracy cost quantified in \refs{sec:resSelection}.

Lower cost per window is only the first requirement, because the pipeline has yet to meet a live stream. A deployment study would connect it to an interrogator through a streaming ingest layer, process many loci concurrently, and measure throughput, end-to-end latency, memory, and energy draw on the embedded platform, the quantities an integrator needs and this thesis does not report. Integration with alarm-management software would then allow a nuisance-alarm budget to be evaluated at natural class prevalence.

\subsection{Evaluate Interpretability with Operators}
\label{sec:chap5FutureOperators}
\autoBookmark{subsec}{Evaluate Interpretability with Operators}
The claim of interpretability rests on named descriptors and attribution scores. Whether these help a security operator decide faster, or trust an alert more appropriately, is an empirical question for a user study that compares alerts with and without feature-level explanations. Such a study would also show which descriptors, or which summaries of them, mean something to the people who act on them.

\section{Concluding Remarks}
\label{sec:chap5Closing}
\autoBookmark{sec}{Concluding Remarks}
This thesis set out to build an accurate, interpretable, and deployable DAS event classifier, and to measure it in a way that others can check. Multi-domain handcrafted features paired with gradient boosting, designed on the Tomasov perimeter-security dataset and applied unchanged to the Cao dataset, outperform the convolutional network published with each dataset and five recent transformer and graph architectures on identical leakage-controlled partitions, with sub-millisecond inference on embedded hardware. That result holds for unseen segments of represented recordings, and it comes with an explicit account of where the artefact falls short: a numerical residue in the normalisation that flatters the Tomasov figures and narrows the lead over the neural baselines, a minority class that only channel-level models recover, a real-time budget set by feature extraction rather than inference, a substantial fall in accuracy when whole recordings or acquisition campaigns are held out, and no transfer of a fitted model between sensing systems.

The more durable contribution concerns how such systems are measured. Two widely used public benchmarks were shown to permit substantial information leakage, and correcting it changed not only the headline figures but the qualitative conclusions. A minority class believed detectable proved unrecoverable from pooled descriptors by any of seven classifiers; a balancing strategy believed superior proved hard to distinguish from its alternatives; a feature domain believed to encode event physics proved to encode acquisition identity as well; and the ordering of published architectures did not survive evaluation under one protocol. None of these corrections required a new algorithm. They required an evaluation protocol that could be verified, and statistics computed over genuinely independent units. Turned on the thesis's own feature code, the same discipline found the residue that qualifies its headline.

Group-aware, cross-campaign evaluation should therefore become standard practice in DAS event classification. Without it, published accuracies cannot be compared with one another and do not predict what a deployed system will do. The partition manifests, verification code, protocol-attribution design, and cross-dataset label mapping released with this thesis are intended to make that practice straightforward to adopt on the two public benchmarks examined here.
